\documentclass[12pt]{article}
\usepackage[margin=1in]{geometry}
\usepackage{hyperref}
\usepackage{enumitem}
\usepackage{parskip}

% Define the framework name locally
\newcommand{\sandoval}{\textsc{Multi-SaLLM}}

\title{\textbf{Response to Reviewers}\\[0.5em]
\large ``Multi-SaLLM: A Multilingual Security Assessment of LLM-Generated Code''}
\author{}
\date{}

\begin{document}
\maketitle

We thank all three reviewers for their thorough and constructive feedback. We have
carefully addressed every comment in the revised manuscript. Below, we respond to
each point in the order it was raised. Section and figure references correspond to
the revised submission.
\bigskip
\hrule


%=============================================================================
\section*{Reviewer 1}
%=============================================================================

\begin{enumerate}[leftmargin=*, label=\textbf{R1.\arabic*.}]

% -----------------------------------------------------------------------
\item \textit{``The multilingual translation with 23 different prompting languages does
not bring a substantial value added to the work, as their relevance to the work is not
clear, while the experimental results do not show significant differences among the
different languages (e.g., in Figure 4, the 23 charts are very similar). It would be
more relevant, instead, to have different programming languages other than Python. It
will be more challenging to measure the security of code with typed and memory unsafe
languages, such as C and C++, with various and complex memory errors, and also with
safer programming languages such as Java or Rust.''}

\textbf{Response.}
We thank the reviewer for this important observation. In direct response, we have
extended the framework and dataset to include \textbf{Java} as a second target
programming language. For each of the 125 Python prompts we manually constructed
complementary Java prompts that preserve the original task semantics and security
context, including equivalent insecure code examples and aligned unit-test cases
(Section~3.1.4). The evaluation pipeline, CodeQL static analysis and JUnit-based
dynamic testing, was extended to handle Java outputs end-to-end
(Sections~3.3.1--3.3.2). Results for Java are reported alongside Python throughout
Sections~4.1--4.4, revealing meaningful cross-language differences in post-repair
compilability, pass@k, and security@k. Regarding the similarity of multilingual
charts: our language-family analysis (Section~4.4, Table~4) confirms that
\textit{natural-language} variation is modest, while \textit{programming language}
(Python vs.\ Java) is the dominant performance factor, consistent with the
reviewer's intuition. We acknowledge in the Limitations section (Section~5.4) that
extending to C/C++ or Rust would require different static analysis toolchains and
memory-safety test oracles, which we plan to address in future work.

% -----------------------------------------------------------------------
\item \textit{``A question that is not addressed concerns the functional correctness of
unsecured code (containing vulnerabilities or run-time errors) because in that case the
behavior is not well-defined.''}

\textbf{Response.}
We agree this is an important nuance. Our framework explicitly separates
\textit{functional correctness} (pass@k, measured by unit tests that verify
input/output behavior) from \textit{security correctness} (vulnerable@k /
security@k, measured by CodeQL and security-specific test assertions). A code
sample is considered functionally correct only if it passes the functional test
method; it is considered secure only if the security test method also passes (or
CodeQL detects no vulnerability). A sample may therefore be functionally correct
yet insecure, or insecure yet non-functional. This separation is what produces the
security--functionality trade-off discussed in Section~5.1: models with higher
pass@k (e.g., GPT-4o-Mini) simultaneously exhibit the highest vulnerability rates
under both assessment techniques. The framework design is described in
Section~3.3.2, and the trade-off is analyzed in Section~5.1.

% -----------------------------------------------------------------------
\item \textit{``Concerning the experimentation, study on the scalability of the approach
is limited due to the limited size of the dataset (only 92 code snippets, which are not
representative enough), the fact that unit tests were created manually, as well as the
prompts.''}

\textbf{Response.}
We have extended the dataset from 92 to \textbf{125 prompts}, adding 25 prompts
derived from real-world vulnerabilities sourced from Common Vulnerabilities and
Exposures (CVEs) and vulnerable code patterns mined from GitHub repositories
(Section~3.1.1). This extension improves ecological validity by grounding the
benchmark in realistic security scenarios rather than synthetic examples alone. We
acknowledge that manual prompt and test-case creation is a threat to scalability
and discuss this in the Limitations section (Section~5.4). However, the manual
process is essential for ensuring that each prompt genuinely exercises a
security-relevant code path with verified functional and security
assertions---a requirement that cannot be trivially automated. The full dataset
and all test cases are publicly released to enable community extension.

% -----------------------------------------------------------------------
\item \textit{``The metrics are too simple, especially the security@k metric which is
not giving additional insights.''}

\textbf{Response.}
The simplicity of the metrics is a deliberate design choice analogous to the
widely adopted pass@k metric. Nevertheless, we have considerably expanded the
metric evaluation in the revised paper by: (i)~reporting vulnerable@k and
security@k for $k \in \{1, 3, 5\}$ instead of only $k=1$; (ii)~computing both
metrics under two complementary assessment techniques (CodeQL static analysis and
unit-test-based dynamic analysis) for both Python and Java, yielding four distinct
metric perspectives (Sections~4.3.1--4.3.2); and (iii)~analyzing the interaction
between temperature, $k$, and security outcomes in detail in Sections~4.3 and
5.1. These results reveal non-trivial trends: for instance, security@k collapses
sharply as $k$ and temperature grow, and a model's apparent ``safety'' may be an
artifact of low compilability rather than robust security behavior.

% -----------------------------------------------------------------------
\item \textit{``There is no information about the CWE coverage (number of treated CWEs
vs total of CWEs for Python code).''}

\textbf{Response.}
Our dataset targets the \textbf{Top 25 Most Dangerous Software Weaknesses} as
applicable to Python, supplemented by CWEs with demonstrative examples from the
MITRE CWE list, CodeQL's vulnerability documentation, and Sonar Rules. The prompt
sources and individual counts are detailed in Section~3.1.1. Each prompt is labelled
with a CWE-ID (Section~3.1.2), and a complete prompt-to-CWE mapping is provided in
the released replication package. We have added a clarifying note on CWE coverage
in Section~3.1.1 to make the scope explicit.

\end{enumerate}

\bigskip\hrule\bigskip

%=============================================================================
\section*{Reviewer 2}
%=============================================================================

\begin{enumerate}[leftmargin=*, label=\textbf{R2.\arabic*.}]

% -----------------------------------------------------------------------
\item \textit{``Static evaluation depends solely on CodeQL without cross-validation
using other mainstream static analysis tools (e.g., Bandit), and the dynamic evaluation
test cases are limited to known CWE vulnerabilities without adversarial scenarios.''}

\textbf{Response.}
We complement the CodeQL static analysis with a fully independent
\textbf{test-based dynamic assessment} (Section~3.3.2). For each prompt we manually
created unit tests containing both a functional test method and a security test
method. Vulnerable@k and Security@k are therefore computed independently under both
techniques and reported side-by-side in Sections~4.3.1 (Static-based) and 4.3.2
(Test-based); the Discussion (Section~5.1) contrasts the two assessments explicitly.
Regarding cross-validation with additional static tools: \sandoval{} is designed to
be \textit{configurable}, making it possible to substitute or add analyzers. We
focused on CodeQL because it supports both Python and Java with taint analysis,
enabling consistent cross-language comparison. Adding tools such as Bandit remains
future work. Regarding adversarial scenarios: our security test methods use crafted
inputs specifically designed to trigger the CWE being tested (e.g., malicious
usernames for SSRF, crafted SQL payloads for injection). Full adversarial fuzzing
is an extension we plan to pursue.

% -----------------------------------------------------------------------
\item \textit{``The dataset would be more representative if code snippets from real
GitHub projects were incorporated. The dataset has major limitations, including support
only for Python and very limited real-world snippet coverage (e.g., only 13 valid Stack
Overflow samples), restricting cross-language generalizability and scenario diversity.''}

\textbf{Response.}
We have addressed both sub-points. First, we added \textbf{25 prompts sourced
directly from real-world CVEs and GitHub vulnerability patterns}
(Section~3.1.1), supplementing the original synthetic/CWE-derived prompts with
ecologically valid scenarios. Second, we have fully extended the benchmark to
include \textbf{Java}: 125 Java prompts with corresponding insecure examples and
unit tests were manually created and evaluated (Sections~3.1.4, 4.1--4.4). We
also note that Stack Overflow snippets represent only one of five sources used in
our dataset; the others (CWE Top 25, CodeQL documentation, Sonar Rules, and the
new GitHub/CVE prompts) together contribute the remaining 112 prompts.

% -----------------------------------------------------------------------
\item \textit{``Although multilingual support uses back-translation and BERTScore, the
lack of human evaluation suggests that small-scale manual quality checks should be
added.''}

\textbf{Response.}
We acknowledge this limitation directly in Section~5.4 (Limitations: Translation
Quality). Our translation pipeline generates 10 candidates per prompt per language
and selects the one with the highest back-translation BERTScore (Section~3.1.3),
providing an automated proxy for semantic fidelity. We agree that native-speaker
human evaluation would further strengthen confidence in translation quality and have
identified it as a direction for future work. Performing such validation across all
22 non-English languages for 125 prompts is beyond the scope of the current study,
consistent with comparable multilingual code benchmarks in the literature (Peng
et al., 2024).

% -----------------------------------------------------------------------
\item \textit{``Insufficient scenario coverage due to very few collected code snippets,
restricting cross-language generalizability and scenario diversity.''}

\textbf{Response.}
As noted in R2.2, the 25 newly added GitHub/CVE-derived prompts directly address
scenario diversity. Together with the Java extension, the revised framework now
evaluates 125 prompts across two programming languages and 23 natural languages.
We also discuss extending scenario coverage as future work in Section~6, including
additional CWE classes and more programming languages.

\end{enumerate}

\bigskip\hrule\bigskip

%=============================================================================
\section*{Reviewer 3}
%=============================================================================

\begin{enumerate}[leftmargin=*, label=\textbf{R3.\arabic*.}]

% -----------------------------------------------------------------------
\item \textit{``In the Abstract section, it is necessary to provide a brief summary of
the findings (Experimental results, comparison with baselines, etc.). Currently, when
reading the abstract, it is not clear what the contributions of the work are, as well
as how it is related to state-of-the-art research.''}

\textbf{Response.}
We have revised the abstract to include a concise summary of key empirical findings.
The updated abstract now reports: (i)~that GPT-4o-Mini achieves the highest
functional correctness across both languages but also carries the highest
vulnerability rates; (ii)~that Python consistently outperforms Java in both pass@k
and security@k; and (iii)~that increasing sampling temperature and $k$ improves
correctness but degrades security guarantees. The three main framework contributions
are also stated explicitly.

% -----------------------------------------------------------------------
\item \textit{``To guide the presentation, I think it makes sense to add a motivating
example to illustrate the applicability of Multi-SaLLM in a real-world scenario.''}

\textbf{Response.}
Figure~1 in Section~2.2 illustrates a concrete motivating scenario: a model
generates a functionally correct yet SQL-injection-vulnerable solution (CWE-89),
and the caption explains how \sandoval{} would detect and quantify this risk
through both static (CodeQL taint analysis) and dynamic (security unit test)
assessment. The figure also shows the secure variant using parameterized queries,
making the applicability of the framework immediate.

% -----------------------------------------------------------------------
\item \textit{``There is Section 1.1., but no Section 1.2. I suggest the authors add
another subsection, or remove the heading Section 1.1.''}

\textbf{Response.}
Thank you for catching this structural inconsistency. The revised manuscript adds
a companion subsection~1.2 (``Paper Organization'') that describes the structure of
the remainder of the paper, giving Section~1.1 (``Extended Contribution'') a proper
sibling. The Introduction now contains two balanced subsections and the section
numbering is consistent throughout.

% -----------------------------------------------------------------------
\item \textit{``The Related Work section is too short, and it does not associate well
the submitted work with state of the art research in the domain. I can see that there
are various papers about security testing of AI generated code, and thus I think it is
necessary to distinguish the current paper from existing work. At least, there is the
following paper: `DeVAIC: A tool for security assessment of AI-generated code'
(\url{https://www.sciencedirect.com/science/article/pii/S0950584924001770}), and it
is highly advisable to relate Multi-SaLLM to DeVAIC.''}

\textbf{Response.}
We have substantially expanded Section~2.4 (Related Work) with twelve additional
references covering recent advances in secure code generation benchmarks and
tools. We now discuss and contrast \sandoval{} with DeVAIC, SVEN, SafeCoder,
CodeGuard+, PromSec, CodeLMSec, CodeSecEval, CWEval, SafeGenBench,
SecRepoBench, and SecureAgentBench. We differentiate \sandoval{} from DeVAIC
and related tools along three axes: (i)~multilingual prompt evaluation across 23
natural languages; (ii)~dual programming-language (Python and Java) assessment;
and (iii)~the combination of static and dynamic analysis under unified security
metrics (security@k, vulnerable@k) that jointly account for both functional
correctness and vulnerability detection.

% -----------------------------------------------------------------------
\item \textit{``An important aspect is the delta of this work compared to that of the
workshop paper. I think it is crucial to compare and contrast this in the Results
section.''}

\textbf{Response.}
Section~1.1 (``Extended Contribution'') now provides a detailed, structured
comparison of the current manuscript versus the original workshop paper, covering:
(i)~the multilingual dataset (22 new languages); (ii)~the multi-programming
extension to Java; (iii)~the 25 real-world GitHub/CVE prompts; (iv)~the expanded
model study (four models across six temperatures); and (v)~the extended metric
analysis ($k \in \{1,3,5\}$ under both static and dynamic assessment for two
programming languages). The Results section (Section~4) quantifies these
extensions through separate Java and Python figures and tables.

% -----------------------------------------------------------------------
\item \textit{``AUSE imposes the author-based citation style (not numbers), please
consult a published paper for more details.''}

\textbf{Response.}
We have updated the manuscript to use the \texttt{sn-mathphys-ay} document class,
which enforces the author-year citation style required by AUSE. All in-text
citations now follow the author-year convention and the reference list is formatted
accordingly.

% -----------------------------------------------------------------------
\item \textit{``Figures 4--9 are pretty small and unreadable, please consider enlarging
them all.''}

\textbf{Response.}
All figures have been enlarged to full column width. Furthermore, the previously
combined multi-language comparison figures have been split into separate Python and
Java variants (appearing in Sections~4.3.1--4.3.2), giving each figure more space
and improving legibility of axis labels, legend entries, and data points.

% -----------------------------------------------------------------------
\item \textit{``Page 12: `When an LLM generate code from [...]' $\rightarrow$
`When an LLM generates code from [...]'''}

\textbf{Response.}
Corrected. The grammatical error has been fixed in the revised manuscript.

\end{enumerate}

\bigskip\hrule\bigskip

\noindent We hope the revised manuscript and this response letter satisfactorily
address all concerns raised. We remain available for any further clarifications.

\bigskip
\noindent\textit{The Authors}

\end{document}
